% Kompiuterijos katedros šablonas
% Template of Department of Computer Science II
% Versija 1.0 2015 m. kovas [ March, 2015]

\documentclass[a4paper,12pt,fleqn]{article}
\usepackage[unicode,colorlinks=false]{hyperref}


\usepackage[utf8x]{inputenc}
%

\usepackage[L7x]{fontenc}
\usepackage{times}
\usepackage{ucs}

 %package to switch the language
\usepackage{etoolbox}

  %set up of the page margins
\usepackage[top=2cm, bottom=2cm, left=3cm, right=1.5cm]{geometry}

 %1.1 line spacing
\linespread{1.1}


  %page numbering at the right side
\usepackage{fancyhdr}
\pagestyle{fancyplain}
\fancyhf{}
\renewcommand{\headrulewidth}{0pt} 
\fancyhfoffset[RO]{0cm}

  %to number at the bottom (exchange lines to number at the top)
\rfoot{\thepage}
  %\rhead{\thepage} %

% \usepackage[usenames,dvipsnames]{pstricks}
\urlstyle{same}
\hypersetup{
%  citecolor=Blue,
%  linkcolor=Blue,
%  urlcolor=Blue
pdfborder={0 0 0 }
}

 %for includegraphics
\usepackage{graphicx}



\usepackage[toc,page]{appendix}


\usepackage{caption}

 %for source codes
\usepackage{listings}
\lstset{commentstyle=\color{red},xleftmargin=10pt, framexleftmargin=6pt, numbersep=1mm, frame=single, numbers=left,numberstyle=\footnotesize,extendedchars=\true, inputencoding=utf8x,basicstyle=\footnotesize,extendedchars=true,
 keywordstyle=\color{black}\bfseries, breaklines=true, breakautoindent=true,framesep=8pt,linewidth=0.95\textwidth
}

 %for algorithms
\usepackage{algorithm}
\usepackage{algorithmic}
 %instead of the above two packages we can use algorithms2e
 %\usepackage[boxed,linesnumbered,vlined,slide]{algorithm2e}

 %special symbols
\usepackage{amsfonts}
\usepackage{amssymb}
\usepackage{amsmath}

 %for theorem like environments
\usepackage{amsthm}

 \usepackage{datetime}
 \renewcommand{\dateseparator}{--}


% SI system units
\usepackage{siunitx}
\sisetup{detect-all}
% Problem with fonts \SI{x.xx}{\micro\metre}, solved with updmap-sys --enable Map=utm.map
\renewcommand{\sfdefault}{uhv}
\renewcommand{\rmdefault}{utm}
\renewcommand{\ttdefault}{ucr}

% List management (itemize, etc.)
\usepackage{enumitem}

\newcommand*{\urlw}[1]{\href{#1}%
            {\nolinkurl{#1}}}

\numberwithin{equation}{section}


\newtoggle{inLithuanian}
 %If the report is in Lithuanian, it is set to true; otherwise, change to false
\settoggle{inLithuanian}{true}

%create file preface.tex for the preface text
%if preface is needed set to true
\newtoggle{needPreface}
\settoggle{needPreface}{false}

\newtoggle{signaturesOnTitlePage}
\settoggle{signaturesOnTitlePage}{true}



\theoremstyle{definition}
\newtheorem{definition}{\keyWordDefinition}
\newtheorem{example}{\keyWordExample}
\def\QED{\unskip\nobreak\hfill\kern5pt$\Box$}

\iftoggle{inLithuanian}{
%\usepackage[L7x]{fontenc}
\usepackage[english,lithuanian]{babel}

\newcommand{\todayiso}{\the\year \dateseparator \twodigit\month \dateseparator \twodigit\day}


\renewcommand{\today}{\number\year\space m. \space \ifcase\month\or
  sausio\or vasario\or kovo\or balandžio\or gegužės\or birželio\or
  liepos\or rugpjūčio\or rugsėjo\or spalio\or lapkričio\or
  gruodžio\fi
  \space\number\day\space d.}


 \usepackage{tocloft}
 \renewcommand\cftsecaftersnum{.} 
 \renewcommand\cftsubsecaftersnum{.} 
 \renewcommand\cftsubsubsecaftersnum{.}

 \usepackage{VUMIFKK}

 \DeclareCaptionLabelFormat{captionlt}{#2 #1}
   %smth is not fine with algorithms 
 \DeclareCaptionLabelFormat{captionltalg}{#2 #1 algoritmas}

 \usepackage{indentfirst}
 \renewcommand{\appendixtocname}{Priedai}
 \renewcommand{\appendixpagename}{Priedai}
 \renewcommand{\contentsname}{Turinys} 

 \renewcommand{\lstlistingname}{išeities kodas}
 \renewcommand{\figurename}{pav}
 \renewcommand{\tablename}{lentelė}


 \captionsetup*[lstlisting]{   
 labelsep=period,labelformat=captionlt
 }
 \captionsetup*[figure]{   
% labelsep=period,
 labelsep=space, %babel redefines pav to pav.
 labelformat=captionlt
 }
 \captionsetup*[table]{   
  labelsep=period,
  labelformat=captionlt
 }
 \renewcommand{\algorithmicrequire}{\textbf{Įvestis:}}
 \renewcommand{\algorithmicensure}{\textbf{Išvestis:}}

 \captionsetup*[algorithm]{   
 labelsep=period,labelformat=captionltalg
 }

\renewcommand{\thmhead}[3]{#2 #1#3}

}
{
%\usepackage[OT1,T1]{fontenc}
%\usepackage[L7x]{fontenc}

\usepackage[english]{babel}
\newcommand{\todayiso}{\twodigit\month \dateseparator \twodigit\day \dateseparator \the\year}
 \captionsetup*[algorithm]{   
 labelsep=period
 }
\captionsetup*[lstlisting]{   
 labelsep=period
 }
 \captionsetup*[figure]{   
 labelsep=period
 }
 \captionsetup*[table]{   
 labelsep=period
 }


}

%some kywords
 \def\keywordAbstract{\iftoggle{inLithuanian}{Santrauka}{Abstract}}
 \def\keywordAbstractOther{\iftoggle{inLithuanian}{Summary}{Santrauka}}
 \def\keyWordIntroduction{\iftoggle{inLithuanian}{Įvadas}{Introduction}}
 \def\keyWordConclusions{\iftoggle{inLithuanian}{Išvados ir rekomendacijos}{Conclusions and Recommendations}}

 \def\keyWordPreface{\iftoggle{inLithuanian}{Pratarmė}{Preface}}
 \def\keyWordAppendice{\iftoggle{inLithuanian}{Priedas}{Appendix}}
 \def\keyWordSignature{\iftoggle{inLithuanian}{parašas}{signature}}
 \def\keyWordDefinition{\iftoggle{inLithuanian}{apibrėžimas}{Definition}}
 \def\keyWordExample{\iftoggle{inLithuanian}{pavyzdys}{Example}}

\newcommand{\bothabstracts}[3]{
\setcounter{secnumdepth}{0}
\newpage
\hspace{2cm}
{\centering{\section{\keywordAbstract}}}

#1
\newpage
\hspace{2cm}
{\centering \section{\keywordAbstractOther}}

\begin{center}{\textbf{#2} }\end{center}

 #3
\setcounter{secnumdepth}{3}
}

 %non-numbered sections: #1 param: for labeling sec:#1, #2 -section title
\newcommand{\sectionWithoutNumber}[2]{\newpage
%\hspace{2cm}
\section*{#1}
\label{sec:#2}
\addcontentsline{toc}{section}{\nameref{sec:#2}}%{#3}
 }



\newcommand{\referenceSources}[1]{
\newpage
\cleardoublepage
\phantomsection
\iftoggle{inLithuanian}{
 \renewcommand{\refname}{Literatūros šaltiniai}

 \addcontentsline{toc}{section}{Literatūros šaltiniai}
 \markboth{\refname}{Literatūros šaltiniai}
 }
{

\addcontentsline{toc}{section}{References}
\markboth{References}{References}
}

\bibliographystyle{plain}
\bibliography{#1}
}



 \newcommand\authorsignature[1]{
\begin{flushright}
 \begin{minipage}[b]{0.45\textwidth}
  \centering
  \rule{\textwidth}{0.5pt}\\
   #1
  \end{minipage}
\end{flushright}
 }




 \newcommand\authorsignatures[5]{%
   \vspace{1cm}
   \authorsignature{#1}
   \ifstrequal{#2}{}{}{\vspace{0.3cm}
     \authorsignature{#2}
     \ifstrequal{#3}{}{}{\vspace{0.3cm}
      \authorsignature{#3}
      \ifstrequal{#4}{}{}{\vspace{0.3cm}
        \authorsignature{#4}
        \ifstrequal{#5}{}{}{\vspace{0.3cm}
         \authorsignature{#5}       
        }
      }
    }
} 
}

\newcommand{\authortitle}{
\iftoggle{signaturesOnTitlePage}{
\tiny{\keyWordSignature}
}{}
}

\newcommand{\depttitlepage}[8]
{
\thispagestyle{empty}
\begin{center}

\includegraphics[width=2cm]{jb_VU_zenklas}

%\vspace{-1cm}

\iftoggle{inLithuanian}
{ 
  VILNIAUS UNIVERSITETAS\\
  MATEMATIKOS IR INFORMATIKOS FAKULTETAS\\
  KOMPIUTERIJOS KATEDRA
}
{
  VILNIUS UNIVERSITY \\
  FACULTY OF MATHEMATICS AND INFORMATICS \\
  DEPARTMENT OF COMPUTER SCIENCE II
}

\vspace{5cm}

#1\\
\vspace{0.5cm}
\textbf{\Large #2}
\end{center}

\vspace{5cm}


\hspace{0.5\textwidth}
\begin{minipage}{0.4\textwidth}
 \begin{flushleft} 
\iftoggle{inLithuanian}
{
 \ifstrequal{#3}{}{}{Atliko:\\[5pt]}
}
{
\ifstrequal{#3}{}{}{Done by:\\[5pt]}
}

%\noindent
\begin{tabular}{@{}lr}%\setlength\tabcolsep{0pt}
\ifstrequal{#3}{}{}{#3&\hspace{2cm}\authortitle\\[5pt]}
\ifstrequal{#4}{}{}{#4&\authortitle\\[5pt]}
\ifstrequal{#5}{}{}{#5&\authortitle\\[5pt]}
\ifstrequal{#6}{}{}{#6&\authortitle\\[5pt]}
\ifstrequal{#7}{}{}{#7&\authortitle\\}
\end{tabular}

\end{flushleft}

\end{minipage}

\vspace{0.5cm}
\hspace{0.5\textwidth}
\begin{minipage}{0.4\textwidth}
 \begin{flushleft} 

\ifstrequal{#8}{}{}
{

\iftoggle{inLithuanian}
{
Vadovas:
}
{
Supervisor:
}

#8

}

\end{flushleft}

\end{minipage}


\vfill

\begin{center}
Vilnius\\
\the\year
\end{center}

\iftoggle{needPreface}{
 \sectionWithoutNumber{\keyWordPreface}{preface}
Pratarmės (Preface) informacija


\iftoggle{inLithuanian}
{
\vspace{\baselineskip}\hfill
\today
}
{
 \vspace{\baselineskip}\hfill \today
}

 \vspace{5cm}

\iftoggle{signaturesOnTitlePage}{}
{
\authorsignatures{#3}{#4}{#5}{#6}{#7}
}
}{}
\newpage
}


\begin{document}
 % #1 -report type, #2 - title, #3-7 students, #8 - supervisor
 \depttitlepage{Darbo/ataskaitos tipas}{Kibernetinio saugumo pažeidžiamos paslaugos ir atakos vektoriaus įgyvendinimo metodų tyrimas}{Deividas Slauzgalvis}{Mindaugas Strakšys} 
 {}{}{}{}% students 2-5
 {dr. Vadovas Vadovaitis}

\tableofcontents


%keywords and notations if needed
\sectionWithoutNumber{Sutartinis terminų žodynas}{keywords}{Pateikiamas terminų sąrašas (jei reikia)}

 %both abstracts
\bothabstracts{Santraukos tekstas rašto darbo kalba...}%tex-file of abstract in original language
{Darbo pavadinimas kita kalba} %if work is in LT this title should be in English
{This is a summary in English...}%tex-file of abstract in other language


 %Introduction section: label is sec:intro
\sectionWithoutNumber{\keyWordIntroduction}{intro}
Įvado tekstas ...



 %the main part
\newpage
\section{Kibernetinio saugumo analizė}
\label{sec:motivation}
\subsection{Kibernetinio saugumo apžvalga}


\subsubsection{Kibernetinio saugumo istorija}
Kibernetinio saugumo industrija, kuri 2015 metais buvo verta 75 milijardų dolerių, prasidėjo 1988 metais, kai Robert T.Morris paleido savo replikuojantį "kirminą" ARPANET tinkle(interneto pradininkas). Šis "kirminas" buvo dalis projekto, kuris skaičiavo interneto dydį užkrėsdamas UNIX operacines sistemas tam kad suskaičiuotų jose esančių prisijungimų prie interneto kiekį. Dėl programavimo klaidos, "kirminas" pradėjo jungtis į tas pačias mašinas daugelį kartų, taip visiškai "užkišdamas" tinklus ir priverdamas sistemas "užlūžti", taip tapdamas pirmuoju tokio tipo įrankiu susilaukusiu didžiulo žiniasklaidos dėmesio. Jo kūrėjas buvo išmestas iš universiteto, nuteistas lygtinai trejiems metams bei nubaustas 10 000 dolerių bauda. \newline Paskutiniame 20a. dešimtmetyje virusai pradėjo smarkiai plisti internetinėje erdvėjė, turbūt tada populiaurisi jų - \textbf{ILOVEYOU} ir \textbf{The Melissa} virusai. Jie užkrėtė dešimtis milijonų kompiuterių visame pasaulyje, privertė žlugti daugelį elektroninio pašto sistemų. Šie virusai pasižymėjo vienu įdomiu dalyku - jie neturėjo apibrėžto tikslo ir taikinio, ir nesiekė jokios finansinės ar politinės naudos. Vienintelis jų tikslas buvo sukelti paniką ir chaosą virtualioje erdvėje. Tam, kad pasipriešinti šiems virusams buvo pradėtos kurtis pirmosios sistemos, skirtos apsiginti nuo virusų (antivirusinės). Taipogi, kompanijos pradėjo šviesti ir edukuoti savo darbuotojus kibernetinio saugumo temomis. Jie buvo išmokyti neatidarinėti neaiškių elektroninių laiškų, tam, kad išvengti "phishing" tipo atakų. Būtent šiuo metu kompanijos rimtai susirūpino kibernetinio saugumo klausimais ir apie tai pradėta diskutuoti viešumoje. \newline 21a. pirmajame dešimtmetyje kibernetinės atakos pasidarė dar rimtesnės ir pavojingesnės. Atakuoti buvo pradėtos kreditinės kortelės ir elektroninės banko sąskaitos. Tarp 2005 ir 2007 metų Albert Gonzalez su savo kibernetinių nusikaltėlių gauja sugebėjo pavogti informaciją iš bent 45,7 milijono kreditinių kortelių. TJX kompanija dėl to patyrė 256 milijonų dolerių nuostolį. Tai buvo puiki pamoka tiek TJX, tiek kitoms kompanijoms. Į kibernetinį saugumą reikėjo pradėti žiūrėti daug rimčiau nei ankščiau, nes tai praleidus pro akis nuostoliai gali būti milžiniški ir pasibaigti kompanijos bankrotu. \newline 2014 metais kibernetiniai nusikaltėliai perėmė milžiniškus kiekius Sony bei Target kompanijų duomenų. Šįkart buvo taikytasį ne tik į finansinę naudą, bet ir į kompanijos reputaciją ir jos darbuotojų gerovę. Tačiau Target kompaniją dėl to prarado apie 162 milijonus dolerių ir daugelį klientų - nes nusikaltėliai turėjo prieigą prie 70 milijonų Target kompanijos klientų asmeninių duomenų. Taipogi nustatyta,kad už Sony kompanijos duomenų kompromitavimą atsakingą Šiaurės Korėja. Taigi, kibernetinės atakos tapo ne tik finansinės naudos siekimo įrankiu, tačiau ir stipriu politiniu(teroristiniu) įrankiu, skirtu aiškintis nesutarimus tarp šalių. Kuo technologijos modernesnės, tuo daugiau naudingos ir svarbios informacijos ten talpinama. Dabar kibernetinė sauga yra vienas svarbiausių prioritėtų kompanijose. Ir tuo rūpintis reikia, kol atakos dar neįvyko (išankstinė atakų prevencija), nes atakos eksplotavimo metu gintis jau per vėlu ir nuostaliai gali būti tragiški.

\subsubsection{Kompiuterinių sistemų atakavimo apžvalga}
Kiekviena kompiuterinė sistema turi trukumų, kuriuos anksčiau ar vėliau atranda programuotojai. Tačiau trūkumo ieškojimas nėra paprastas procesas, jis susideda iš daugybės etapų, kurie kas kart gali skirtis dėl skirtingų atakų tipų. Tačiau kiekvienos atakos pradinis taškas, nuo kurio prasideda atakos planavimas, yra pasiruošimo stadija.

Kibernetinių atakų pasiruošimo stadijoje svarbiausias žingsnis yra susipažinimas su sistema, kurią ruošiamsi atakuoti. Šiame etape svarbu sužinoti sistemos aukos architektūra, naudojamas programas bei kokių technologijų pagalba ši sistema egzistuoja internetinėje erdvėje.
\subsubsection{Kompiuterinių sistemų saugumo apžvalga}
Hakeriai šiais laikais dažniausiai siejami kaip blogiukai, kurie įsilaužineja į sistemas ar kitaip stengiasi pakenkti kompiuterinėms sistemoms. Tačiau yra hakerių, kurie ieško trūkumų sistemose tam, kad būtų apsisaugota nuo kenkėjiškų hakerių.
\subsection{Pagrindiniai kibernetinių atakų tipai}
\subsubsection{SQL injekcijos}
\subsubsection{Phishingas}
\subsubsection{Kenkėjiškos programos}
\subsubsection{Slaptažodžių laužimas}
\subsubsection{DDOS/DOS atakos}
\section{Pasirinktų atakos vektorių analizė}

\section{Pasirinktų atakos vektorių implementacija}

\section{Šaltiniai}
 %Conclusions section
\sectionWithoutNumber{\keyWordConclusions}{conclusion}
Išvados bei rekomendacijos.


%ateities darbų gairės, planas/next steps of the work
\sectionWithoutNumber{Ateities tyrimų planas}{future}{Pristatomi ateities darbai ir/ar jų planas, gairės tolimesniems darbams....}

 %file literatureSources.bib
\referenceSources{literatureSources}



%% this part is optional
\newpage
\begin{appendices}
Dokumentą sudaro du priedai: \ref{app:a} priede  ....
\newpage
\section{Pirmojo priedo pavadinimas}
\label{app:a}
Pirmojo priedo tekstas ...

\newpage
\section{Antrojo priedo pavadinimas}
Antrojo priedo tekstas ...

\end{appendices}


\end{document}
